\chapter{Implementation}
\label{chapter4}

\section{Implementation Overview}
This chapter details the implementation of the authorization mechanism. The implementation spans
three layers:
(i) ISA extension and decode/execute integration for the \texttt{kc} instruction, (ii) a
verification datapath with protected key storage and an Authenticator FSM, and (iii) host
toolchain and application support for compilation, inspection, authorization exchange, and
programming.

\begin{table}[H]
  \centering
  \caption{Implementation components}
  \label{tab:impl-components}
  \begin{tabular}{@{}p{3.2cm}p{8.8cm}@{}}
    \toprule
    \textbf{Component} & \textbf{Implementation details} \\
    \midrule
    ISA extension & Custom \texttt{kc} opcode and decode path; I-type format \\
    Verification datapath & Constant-time comparator; protected NVM access \\
    Authenticator FSM & \textsf{LOCKED}/\textsf{VERIFY}/\textsf{AUTHORIZED}/\textsf{HALT} \\
    Debug gating & \texttt{debug\_enable} released only on match \\
    Toolchain & GCC/binutils support for emitting \texttt{kc} \\
    Host application & Compile, hex view, key exchange, connect, program \\
    \bottomrule
  \end{tabular}
\end{table}

\section{ISA Extension and Firmware Support}
The custom instruction \texttt{kc} is introduced as an I-type instruction with an immediate
field that conveys a token fragment. The decode/execute path recognizes the opcode pattern and
forwards operands and context to the verification datapath. The instruction is emitted by the
toolchain and preserved through compilation and ELF$\rightarrow$HEX conversion, enabling
inspection and audit in build artifacts.

\section{Protected Key Storage and Verification}
Authorization keys are provisioned into protected non-volatile memory that is not mapped into
the normal load/store address space. Verification is performed by a constant-time comparator,
which outputs a one-bit decision \texttt{match} without exposing the reference value. The
Authenticator FSM consumes this decision to gate execution and debug enablement.

\section{Host Workflow and Authorization Exchange}
The host workflow consists of compilation, hex inspection, authorization exchange over UART,
device connection, and image programming. Authorization uses a short, explicit sequence
(Start, Key-1..Key-3, Enable, End). The host application logs programming events and key
match results, supporting traceability and audit.

\section{Summary}
The implementation integrates a custom ISA instruction, protected key storage, a constant-time
verifier, and an Authenticator FSM to gate execution and debug access. Toolchain support ensures
that the authorization primitive is emitted and preserved across the build pipeline, while the
host workflow provides compilation, inspection, authorization exchange, and programming steps.
Evaluation results and evidence artifacts are presented in Chapter~5.
